\documentclass[11pt,a4paper,modern]{FFExercises}
\usepackage[english,greek]{babel}
\usepackage[utf8]{inputenc}
\usepackage{nimbusserif}
\usepackage[T1]{fontenc}
\usepackage{amsmath}
\let\myBbbk\Bbbk
\let\Bbbk\relax
\usepackage[amsbb,subscriptcorrection,zswash,mtpcal,mtphrb,mtpfrak]{mtpro2}
\usepackage{graphicx,multicol,multirow,enumitem,tabularx,mathimatika,gensymb,venndiagram,hhline,longtable,tkz-euclide,fontawesome5,eurosym,tcolorbox,tabularray,tikzpagenodes,relsize,siunitx,eurosym}
\definecolor{xrwma}{HTML}{aa1212}
\usetikzlibrary{calc}
\usetikzlibrary{positioning}
\tcbuselibrary{skins,theorems,breakable}
\renewcommand{\textstigma}{\textsigma\texttau}
\renewcommand{\textdexiakeraia}{}
\usepackage{svg}
\ekthetesdeiktes

\begin{document}

\titlos{Μαθηματικά Προσανατολισμού}{Γ΄ Λυκείου}{Λύσεις - Ακρότατα Συνάρτησης}

%=============================================
% ΟΜΑΔΑ Α: ΛΥΣΕΙΣ
%=============================================
\paragraph{Ομάδα Α: Λύσεις}

\askhsh % 1
\textbf{Λύση}
\begin{alist}
    \item $A = \mathbb{R}$. Η $f$ είναι πολυωνυμική, άρα συνεχής στο $\mathbb{R}$.
    \item $f'(x) = 3x^2 - 6x = 3x(x-2)$.
    \item $f'(x) = 0 \Leftrightarrow x = 0$ ή $x = 2$. Μονοτονία: $\nearrow (-\infty, 0]$, $\searrow [0, 2]$, $\nearrow [2, +\infty)$.
    \item Τοπικό μέγιστο $f(0) = 5$, τοπικό ελάχιστο $f(2) = 1$.
    \item $f(x) = 3 \Leftrightarrow x^3 - 3x^2 + 2 = 0$. Αφού $f(2) = 1 < 3 < 5 = f(0)$, από Bolzano στο $[0, 2]$ υπάρχει ρίζα.
\end{alist}

\askhsh % 2
\textbf{Λύση}
\begin{alist}
    \item $f'(x) = -6x^2 + 18x - 12 = -6(x^2 - 3x + 2) = -6(x-1)(x-2)$.
    \item $f'(x) = 0 \Leftrightarrow x = 1$ ή $x = 2$.
    \item Πρόσημα: $f' < 0$ στο $(-\infty, 1)$, $f' > 0$ στο $(1, 2)$, $f' < 0$ στο $(2, +\infty)$.
    \item Τοπικό ελάχιστο $f(1) = -1$, τοπικό μέγιστο $f(2) = 0$.
    \item Πίνακας προσήμου-μονοτονίας: $- \searrow 1 \nearrow 2 \searrow +$.
\end{alist}

\askhsh % 3
\textbf{Λύση}
\begin{alist}
    \item $f'(x) = 4x^3 - 12x^2 = 4x^2(x - 3)$.
    \item $f'(x) = 0 \Leftrightarrow x = 0$ (διπλή ρίζα) ή $x = 3$.
    \item Στο $x = 0$ δεν αλλάζει πρόσημο η $f'$ (παραμένει $< 0$), άρα δεν υπάρχει ακρότατο. Στο $x = 3$ αλλάζει από $-$ σε $+$, άρα ελάχιστο.
    \item $\searrow (-\infty, 3]$, $\nearrow [3, +\infty)$.
    \item $\lim_{x \to \pm\infty} f(x) = +\infty$.
\end{alist}

\askhsh % 4
\textbf{Λύση}
\begin{alist}
    \item $f'(x) = 12x^3 - 24x^2 + 12x = 12x(x^2 - 2x + 1) = 12x(x-1)^2$.
    \item $f'(x) = 0 \Leftrightarrow x = 0$ ή $x = 1$.
    \item Στο $x = 1$ (διπλή ρίζα) δεν αλλάζει πρόσημο. Στο $x = 0$ αλλάζει από $-$ σε $+$, άρα ελάχιστο.
    \item $f_{\min} = f(0) = 1$.
\end{alist}

\askhsh % 5
\textbf{Λύση}
\begin{alist}
    \item $f'(x) = 3x^2 - 12x + 9 = 3(x^2 - 4x + 3) = 3(x-1)(x-3)$.
    \item $f'(x) = 0 \Leftrightarrow x = 1$ ή $x = 3$.
    \item Στο $x = 1$: μέγιστο, στο $x = 3$: ελάχιστο. (Ή: $f'(x) = 3(x-1)(x-3) = 3(x-2)^2 - 3$, αλλά αυτό δεν αλλάζει την ανάλυση). Μέγιστο $f(1) = 6$, ελάχιστο $f(3) = 2$.
    \item Η $f$ δεν είναι γνησίως μονότονη αφού έχει δύο διακριτά ακρότατα.
\end{alist}

\askhsh % 6
\textbf{Λύση}
\begin{alist}
    \item $f(-x) = -x^5 + 5x^3 = -(x^5 - 5x^3) = -f(x)$, άρα περιττή.
    \item $f'(x) = 5x^4 - 15x^2 = 5x^2(x^2 - 3)$. Κρίσιμα σημεία: $x = 0, \pm\sqrt{3}$.
    \item Μέγιστο $f(-\sqrt{3}) = 6\sqrt{3}$, ελάχιστο $f(\sqrt{3}) = -6\sqrt{3}$.
    \item $\lim_{x \to +\infty} f(x) = +\infty$, $\lim_{x \to -\infty} f(x) = -\infty$. Δεν υπάρχουν ολικά ακρότατα αφού το σ.τ. είναι $\mathbb{R}$.
    \item $f(x) = 0 \Leftrightarrow x^3(x^2 - 5) = 0 \Leftrightarrow x = 0, \pm\sqrt{5}$. Τρεις ρίζες.
\end{alist}

%=============================================
% ΟΜΑΔΑ Β: ΛΥΣΕΙΣ
%=============================================
\paragraph{Ομάδα Β: Λύσεις}

\askhsh % 7
\textbf{Λύση}
\begin{alist}
    \item $A = \mathbb{R}^*$.
    \item $f'(x) = 1 - \frac{1}{x^2} = \frac{x^2 - 1}{x^2}$.
    \item $f'(x) = 0 \Leftrightarrow x = \pm 1$.
    \item Για $x < -1$: $f' > 0$. Για $-1 < x < 0$: $f' < 0$. Για $0 < x < 1$: $f' < 0$. Για $x > 1$: $f' > 0$.
    \item Τοπικό μέγιστο $f(-1) = -2$, τοπικό ελάχιστο $f(1) = 2$.
\end{alist}

\askhsh % 8
\textbf{Λύση}
\begin{alist}
    \item $A = \mathbb{R}$.
    \item $f'(x) = \frac{1 \cdot (x^2+1) - x \cdot 2x}{(x^2+1)^2} = \frac{1 - x^2}{(x^2+1)^2}$.
    \item $f'(x) = 0 \Leftrightarrow x = \pm 1$.
    \item $f' > 0$ στο $(-1, 1)$, $f' < 0$ αλλού.
    \item Ελάχιστο $f(-1) = -\frac{1}{2}$, μέγιστο $f(1) = \frac{1}{2}$. Σύνολο τιμών $\left[-\frac{1}{2}, \frac{1}{2}\right]$.
\end{alist}

\askhsh % 9
\textbf{Λύση}
\begin{alist}
    \item $A = \mathbb{R}$.
    \item $f'(x) = \frac{(2x-2)(x^2+1) - (x^2-2x+1)(2x)}{(x^2+1)^2} = \frac{2x^3 + 2x - 2x^2 - 2 - 2x^3 + 4x^2 - 2x}{(x^2+1)^2} = \frac{2x^2 - 2}{(x^2+1)^2} = \frac{2(x-1)(x+1)}{(x^2+1)^2}$.
    \item $f'(x) = 0 \Leftrightarrow x = \pm 1$.
    \item Ελάχιστο $f(1) = 0$, μέγιστο $f(-1) = \frac{4}{2} = 2$.
    \item $f(x) \in [0, 2]$.
\end{alist}

\askhsh % 10
\textbf{Λύση}
\begin{alist}
    \item $A = \mathbb{R} \setminus \{-1\}$.
    \item $f'(x) = \frac{2x(x+1) - x^2}{(x+1)^2} = \frac{x^2 + 2x}{(x+1)^2} = \frac{x(x+2)}{(x+1)^2}$.
    \item $f'(x) = 0 \Leftrightarrow x = 0$ ή $x = -2$.
    \item Μέγιστο $f(-2) = \frac{4}{-1} = -4$, ελάχιστο $f(0) = 0$.
    \item (Στο $(-\infty, -2]$: $\nearrow$, $[-2, -1)$: $\searrow$, $(-1, 0]$: $\searrow$, $[0, +\infty)$: $\nearrow$).
\end{alist}

\askhsh % 11
\textbf{Λύση}
\begin{alist}
    \item $A = \mathbb{R}$.
    \item $f'(x) = \frac{2x(x^2+4) - (x^2-4)(2x)}{(x^2+4)^2} = \frac{2x(x^2+4-x^2+4)}{(x^2+4)^2} = \frac{16x}{(x^2+4)^2}$.
    \item $f'(x) = 0 \Leftrightarrow x = 0$.
    \item Ελάχιστο $f(0) = -1$. (Η $f$ φθίνει για $x < 0$ και αυξάνει για $x > 0$).
\end{alist}

\askhsh % 12
\textbf{Λύση}
\begin{alist}
    \item $A = \mathbb{R} \setminus \{\pm 1\}$.
    \item $f'(x) = \frac{2(x^2-1) - 2x \cdot 2x}{(x^2-1)^2} = \frac{-2x^2 - 2}{(x^2-1)^2} = \frac{-2(x^2+1)}{(x^2-1)^2} < 0$ για κάθε $x \in A$.
    \item Δεν υπάρχουν κρίσιμα σημεία.
    \item Η $f$ είναι γνησίως φθίνουσα σε κάθε διάστημα του πεδίου ορισμού.
\end{alist}

%=============================================
% ΟΜΑΔΑ Γ: ΛΥΣΕΙΣ
%=============================================
\paragraph{Ομάδα Γ: Λύσεις}

\askhsh % 13
\textbf{Λύση}
\begin{alist}
    \item $A = \mathbb{R}$.
    \item $f'(x) = e^x + x \cdot e^x = e^x(1 + x)$.
    \item $f'(x) = 0 \Leftrightarrow x = -1$.
    \item Φθίνουσα στο $(-\infty, -1]$, αύξουσα στο $[-1, +\infty)$.
    \item Ελάχιστο $f(-1) = -e^{-1} = -\frac{1}{e}$.
\end{alist}

\askhsh % 14
\textbf{Λύση}
\begin{alist}
    \item $f'(x) = e^{-x} - x \cdot e^{-x} = e^{-x}(1 - x)$.
    \item $f'(x) = 0 \Leftrightarrow x = 1$.
    \item Αύξουσα στο $(-\infty, 1]$, φθίνουσα στο $[1, +\infty)$.
    \item Μέγιστο $f(1) = \frac{1}{e}$.
    \item $\lim_{x \to -\infty} f(x) = -\infty$, $\lim_{x \to +\infty} f(x) = 0$.
\end{alist}

\askhsh % 15
\textbf{Λύση}
\begin{alist}
    \item $f'(x) = \frac{\frac{1}{x} \cdot x - \ln x}{x^2} = \frac{1 - \ln x}{x^2}$.
    \item $f'(x) = 0 \Leftrightarrow \ln x = 1 \Leftrightarrow x = e$.
    \item Αύξουσα στο $(0, e]$, φθίνουσα στο $[e, +\infty)$.
    \item Μέγιστο $f(e) = \frac{1}{e}$.
    \item $\lim_{x \to 0^+} f(x) = \lim_{x \to 0^+} \frac{\ln x}{x} = -\infty$, $\lim_{x \to +\infty} f(x) = 0$ (L'Hospital).
\end{alist}

\askhsh % 16
\textbf{Λύση}
\begin{alist}
    \item $f'(x) = 2x \ln x + x^2 \cdot \frac{1}{x} = 2x \ln x + x = x(2\ln x + 1)$.
    \item $f'(x) = 0 \Leftrightarrow 2\ln x + 1 = 0 \Leftrightarrow x = e^{-1/2} = \frac{1}{\sqrt{e}}$.
    \item Φθίνουσα στο $(0, \frac{1}{\sqrt{e}}]$, αύξουσα στο $[\frac{1}{\sqrt{e}}, +\infty)$.
    \item Ελάχιστο $f\left(\frac{1}{\sqrt{e}}\right) = \frac{1}{e} \cdot \left(-\frac{1}{2}\right) = -\frac{1}{2e}$.
\end{alist}

\askhsh % 17
\textbf{Λύση}
\begin{alist}
    \item $f'(x) = e^x - 1$.
    \item $f'(x) = 0 \Leftrightarrow x = 0$.
    \item Φθίνουσα στο $(-\infty, 0]$, αύξουσα στο $[0, +\infty)$. Ελάχιστο $f(0) = 0$.
    \item Αφού $f(0) = 0$ είναι ελάχιστο, $f(x) \ge 0$ για κάθε $x$.
\end{alist}

\askhsh % 18
\textbf{Λύση}
\begin{alist}
    \item $f'(x) = \frac{1}{x} - 1$.
    \item $f'(x) = 0 \Leftrightarrow x = 1$.
    \item Αύξουσα στο $(0, 1]$, φθίνουσα στο $[1, +\infty)$. Μέγιστο $f(1) = 0$.
    \item Αφού $f(1) = 0$ είναι μέγιστο, $f(x) \le 0$, δηλαδή $\ln x \le x - 1$.
\end{alist}

%=============================================
% ΟΜΑΔΑ Δ: ΛΥΣΕΙΣ
%=============================================
\paragraph{Ομάδα Δ: Λύσεις}

\askhsh % 19
\textbf{Λύση}
\begin{alist}
    \item $f'(x) = 2x - 2\alpha$.
    \item $f'(x) = 0 \Leftrightarrow x = \alpha$.
    \item $\alpha = 2$.
    \item $f(2) = 4 - 8 + 3 = -1$.
\end{alist}

\askhsh % 20
\textbf{Λύση}
\begin{alist}
    \item $f'(x) = 3\alpha x^2 + 2\beta x - 6$.
    \item $f'(1) = 0$: $3\alpha + 2\beta - 6 = 0$. $f'(-2) = 0$: $12\alpha - 4\beta - 6 = 0$. Λύνοντας: $\alpha = 1$, $\beta = \frac{3}{2}$.
    \item $f''(x) = 6x + 3$. $f''(1) = 9 > 0$ (ελάχιστο), $f''(-2) = -9 < 0$ (μέγιστο).
    \item $f(1) = 1 + \frac{3}{2} - 6 + 1 = -\frac{5}{2}$, $f(-2) = -8 + 6 + 12 + 1 = 11$.
\end{alist}

\askhsh % 21
\textbf{Λύση}
\begin{alist}
    \item $f'(x) = 3x^2 - 6\alpha x = 3x(x - 2\alpha)$.
    \item $f'(x) = 0 \Leftrightarrow x = 0$ ή $x = 2\alpha$.
    \item $2\alpha = 2 \Rightarrow \alpha = 1$.
    \item $f'(x) = 3x(x-2)$. $\nearrow(-\infty, 0]$, $\searrow[0, 2]$, $\nearrow[2, +\infty)$.
\end{alist}

\askhsh % 22
\textbf{Λύση}
\begin{alist}
    \item $f'(x) = e^x - \alpha$.
    \item $f'(x) = 0 \Leftrightarrow e^x = \alpha \Leftrightarrow x = \ln \alpha$ (για $\alpha > 0$).
    \item $\ln \alpha = 0 \Rightarrow \alpha = 1$.
    \item $f(0) = 1 - 0 = 1$.
\end{alist}

\askhsh % 23
\textbf{Λύση}
\begin{alist}
    \item $f'(x) = \frac{1}{x} - \lambda$.
    \item $f'(x) = 0 \Leftrightarrow x = \frac{1}{\lambda}$.
    \item $f(1) = \ln 1 - \lambda = -\lambda = -1 \Rightarrow \lambda = 1$. Επίσης $f'(1) = 1 - 1 = 0$ ✓.
    \item $f'(x) = \frac{1-x}{x}$. Για $x < 1$: $f' > 0$ (αύξουσα). Για $x > 1$: $f' < 0$ (φθίνουσα). Άρα μέγιστο.
\end{alist}

\askhsh % 24
\textbf{Λύση}
\begin{alist}
    \item $f'(x) = 3x^2 + 2\alpha x + \beta$.
    \item $f(1) = 0$: $1 + \alpha + \beta - 2 = 0 \Rightarrow \alpha + \beta = 1$. $f'(1) = 0$: $3 + 2\alpha + \beta = 0$. Λύνοντας: $\alpha = -4$, $\beta = 5$.
    \item $f'(x) = 3x^2 - 8x + 5 = (3x-5)(x-1)$. Στο $x=1$ αλλάζει από $+$ σε $-$, άρα μέγιστο.
    \item $f'(x) = 0$ για $x = 1$ και $x = \frac{5}{3}$. Στο $\frac{5}{3}$ ελάχιστο.
\end{alist}

\askhsh % 25
\textbf{Λύση}
\begin{alist}
    \item $f'(x) = 3x^2 + 2\alpha x + 3$.
    \item $\Delta = 4\alpha^2 - 36$.
    \item Δεν έχει ακρότατα όταν $\Delta \le 0$, δηλαδή $|\alpha| \le 3$.
    \item Για $|\alpha| < 3$: $\Delta < 0$, άρα $f'(x) > 0$ πάντα, άρα γνησίως αύξουσα.
\end{alist}

\askhsh % 26
\textbf{Λύση}
\begin{alist}
    \item $f'(x) = \frac{\alpha}{2\sqrt{x}} - 1$.
    \item $f'(x) = 0 \Leftrightarrow \frac{\alpha}{2\sqrt{x}} = 1 \Leftrightarrow x = \frac{\alpha^2}{4}$.
    \item $\frac{\alpha^2}{4} = 4 \Rightarrow \alpha^2 = 16 \Rightarrow \alpha = 4$ (θετικό).
    \item $f(4) = 4 \cdot 2 - 4 = 4$.
\end{alist}

%=============================================
% ΟΜΑΔΑ Ε: ΛΥΣΕΙΣ
%=============================================
\paragraph{Ομάδα Ε: Λύσεις}

\askhsh % 27
\textbf{Λύση}
\begin{alist}
    \item $g(x) = e^x - x - 1$.
    \item $g'(x) = e^x - 1$.
    \item $g'(x) = 0 \Leftrightarrow x = 0$. Ελάχιστο $g(0) = 0$.
    \item $g(x) \ge g(0) = 0 \Rightarrow e^x \ge x + 1$.
\end{alist}

\askhsh % 28
\textbf{Λύση}
\begin{alist}
    \item $g(x) = \ln x - x + 1$.
    \item $g'(x) = \frac{1}{x} - 1$.
    \item $g'(x) = 0 \Leftrightarrow x = 1$. Μέγιστο $g(1) = 0$.
    \item $g(x) \le g(1) = 0 \Rightarrow \ln x \le x - 1$.
\end{alist}

\askhsh % 29
\textbf{Λύση}
\begin{alist}
    \item $f(x) = x^2 - 2\ln x$.
    \item $f'(x) = 2x - \frac{2}{x} = \frac{2(x^2-1)}{x}$.
    \item $f'(x) = 0$ για $x = 1$ (στο $(0, +\infty)$).
    \item $f(1) = 1$. Αφού είναι ελάχιστο, $f(x) \ge 1$.
\end{alist}

\askhsh % 30
\textbf{Λύση}
\begin{alist}
    \item $g(x) = \sin x - x$.
    \item $g'(x) = \cos x - 1$.
    \item $g'(x) \le 0$ για κάθε $x$ (αφού $\cos x \le 1$). Ισότητα μόνο για $x = 2k\pi$.
    \item Η $g$ είναι γνησίως φθίνουσα. Για $x > 0$: $g(x) < g(0) = 0 \Rightarrow \sin x < x$.
\end{alist}

\askhsh % 31
\textbf{Λύση}
\begin{alist}
    \item $g(x) = x - 2\sqrt{x} + 1$.
    \item $g(x) = (\sqrt{x} - 1)^2$.
    \item $(\sqrt{x} - 1)^2 \ge 0 \Rightarrow x - 2\sqrt{x} + 1 \ge 0 \Rightarrow 2\sqrt{x} \le x + 1$.
    \item Ισότητα για $\sqrt{x} = 1$, δηλαδή $x = 1$.
\end{alist}

\askhsh % 32
\textbf{Λύση}
\begin{alist}
    \item $h(x) = x - 1 - \ln x$.
    \item $h'(x) = 1 - \frac{1}{x}$.
    \item $h'(x) = 0 \Leftrightarrow x = 1$. Ελάχιστο $h(1) = 0$.
    \item $h(x) \ge 0 \Rightarrow x - 1 \ge \ln x$.
\end{alist}

%=============================================
% ΟΜΑΔΑ ΣΤ: ΛΥΣΕΙΣ
%=============================================
\paragraph{Ομάδα ΣΤ: Λύσεις}

\askhsh % 33
\textbf{Λύση}
\begin{alist}
    \item $f'(x) = 2x - \frac{2}{x^2} = \frac{2(x^3 - 1)}{x^2}$.
    \item $f'(x) = 0 \Leftrightarrow x = 1$.
    \item Φθίνουσα στο $(0, 1]$, αύξουσα στο $[1, +\infty)$. Ελάχιστο $f(1) = 3$.
    \item $\lim_{x \to 0^+} f(x) = +\infty$, $\lim_{x \to +\infty} f(x) = +\infty$.
    \item $f(x) \in [3, +\infty)$.
\end{alist}

\askhsh % 34
\textbf{Λύση}
\begin{alist}
    \item $f(x) = e^{x \ln x}$.
    \item $f'(x) = e^{x \ln x} \cdot (\ln x + 1) = x^x(\ln x + 1)$.
    \item $f'(x) = 0 \Leftrightarrow \ln x = -1 \Leftrightarrow x = \frac{1}{e}$.
    \item $f\left(\frac{1}{e}\right) = \left(\frac{1}{e}\right)^{1/e} = e^{-1/e}$.
\end{alist}

\askhsh % 35
\textbf{Λύση}
\begin{alist}
    \item $f'(x) = 3x^2 - 3$.
    \item $f'(x) = 0 \Leftrightarrow x = \pm 1$. Μέγιστο στο $-1$, ελάχιστο στο $1$.
    \item $f(-1) = -1 + 3 + 5 = 7$, $f(1) = 1 - 3 + 5 = 3$.
    \item Αφού $f(-1) = 7 > 0$ και $f(1) = 3 > 0$ και $\lim_{x \to -\infty} f(x) = -\infty$, μία ρίζα στο $(-\infty, -1)$.
\end{alist}

\askhsh % 36
\textbf{Λύση}
\begin{alist}
    \item $f'(x) = 3(x-1)(x+1)$. Μέγιστο στο $-1$, ελάχιστο στο $1$.
    \item $f(-1) = 3$, $f(1) = -1$.
    \item Επειδή $f(-1) = 3 > 0$ και $f(1) = -1 < 0$ και όρια $\pm\infty$: 3 ρίζες.
    \item Μία στο $(-\infty, -1)$, μία στο $(-1, 1)$, μία στο $(1, +\infty)$.
\end{alist}

\askhsh % 37
\textbf{Λύση}
\begin{alist}
    \item $f'(x) = \cos x - \sin x$.
    \item $f'(x) = 0 \Leftrightarrow \tan x = 1 \Leftrightarrow x = \frac{\pi}{4}, \frac{5\pi}{4}$.
    \item $f(0) = 1$, $f(\frac{\pi}{4}) = \sqrt{2}$, $f(\frac{5\pi}{4}) = -\sqrt{2}$, $f(2\pi) = 1$. Μέγιστο $\sqrt{2}$, ελάχιστο $-\sqrt{2}$.
    \item $f(x) \in [-\sqrt{2}, \sqrt{2}]$.
\end{alist}

\askhsh % 38
\textbf{Λύση}
\begin{alist}
    \item $g(x) = e^x - x$. $g'(x) = e^x - 1$. Ελάχιστο $g(0) = 1 > 0$. Δεν τέμνονται.
    \item $h(x) = \ln x - x + 2$. $h'(x) = \frac{1}{x} - 1$. Μέγιστο $h(1) = 1 > 0$. Με Bolzano: $h(e) = 1 - e + 2 = 3 - e \approx 0{,}28 > 0$, $h(4) = \ln 4 - 2 \approx -0{,}61 < 0$. Μία ρίζα.
    \item $f(-1) = 6 > 0$, $f(1) = 2 > 0$, ελάχιστο $f(1) = 2 > 0$. Για $x \to -\infty$: $f \to -\infty$. Μία ρίζα.
    \item $g(x) = x^4 - 4x + 1$. $g'(x) = 4x^3 - 4 = 4(x-1)(x^2+x+1)$. Ελάχιστο $g(1) = -2 < 0$. Όρια $+\infty$. 2 ρίζες.
\end{alist}

%=============================================
% ΟΜΑΔΑ Ζ: ΛΥΣΕΙΣ
%=============================================
\paragraph{Ομάδα Ζ: Λύσεις}

\askhsh % 39
\textbf{Λύση}
\begin{alist}
    \item $2x + 2y = 40 \Rightarrow y = 20 - x$.
    \item $E(x) = x(20 - x) = 20x - x^2$.
    \item $A = (0, 20)$.
    \item $E'(x) = 20 - 2x = 0 \Rightarrow x = 10$. Άρα $y = 10$.
    \item $E_{\max} = 100$ cm$^2$.
\end{alist}

\askhsh % 40
\textbf{Λύση}
\begin{alist}
    \item $y = 10 - x$.
    \item $P(x) = x(10 - x)$.
    \item $A = (0, 10)$.
    \item $P'(x) = 10 - 2x = 0 \Rightarrow x = 5$. Άρα $x = y = 5$.
\end{alist}

\askhsh % 41 (αν χρειάζεται - αλλά έχουμε 40)
\textbf{Λύση (Κουτί)}
\begin{alist}
    \item $x^2 h = 32 \Rightarrow h = \frac{32}{x^2}$.
    \item $E(x) = x^2 + 4xh = x^2 + \frac{128}{x}$.
    \item $E'(x) = 2x - \frac{128}{x^2} = 0 \Rightarrow x^3 = 64 \Rightarrow x = 4$. $h = 2$.
    \item $E_{\min} = 16 + 32 = 48$ m$^2$.
\end{alist}

\askhsh
\textbf{Λύση (Απόσταση από ευθεία)}
\begin{alist}
    \item $d = \sqrt{x^2 + (2x+3)^2} = \sqrt{5x^2 + 12x + 9}$.
    \item $D(x) = 5x^2 + 12x + 9$. $D'(x) = 10x + 12 = 0 \Rightarrow x = -\frac{6}{5}$.
    \item $x = -1{,}2$.
    \item $M\left(-\frac{6}{5}, \frac{3}{5}\right)$.
\end{alist}

\askhsh
\textbf{Λύση (Εφαπτομένη)}
\begin{alist}
    \item $f'(x_0) = -2x_0 + 4$. Εφαπτομένη: $y = (-2x_0+4)(x-x_0) + f(x_0)$.
    \item Τομή με $y=0$: ... Τομή με $x=0$: ...
    \item Εμβαδόν $E(x_0) = \frac{1}{2}|x_y||y_x|$.
    \item Ελαχιστοποίηση...
\end{alist}

\askhsh
\textbf{Λύση (Απόσταση από $y=\ln x$)}
\begin{alist}
    \item $d^2 = x^2 + \ln^2 x$.
    \item $(d^2)' = 2x + \frac{2\ln x}{x} = \frac{2(x^2 + \ln x)}{x}$. Μελέτη της $g(x) = x^2 + \ln x$.
    \item Αριθμητικά: ρίζα κοντά στο $0{,}65$.
    \item Η μέθοδος είναι ελαχιστοποίηση του τετραγώνου της απόστασης.
\end{alist}

%=============================================
% ΟΜΑΔΑ Η: ΛΥΣΕΙΣ
%=============================================
\paragraph{Ομάδα Η: Λύσεις}

\askhsh
\begin{alist}
    \item \textbf{Λάθος} (π.χ. $f(x) = x^3$ στο $x=0$).
    \item \textbf{Λάθος} (μπορεί να μην είναι παραγωγίσιμη, π.χ. $f(x) = -|x|$).
    \item \textbf{Σωστό}.
    \item \textbf{Λάθος} (μπορεί να είναι στα άκρα ή σε σημεία μη παραγωγισιμότητας).
\end{alist}

\askhsh
\begin{alist}
    \item \textbf{Σωστό} (Θεώρημα Μέγιστης-Ελάχιστης Τιμής).
    \item \textbf{Σωστό} (αν η $f$ ορίζεται και είναι συνεχής στο $\mathbb{R}$).
    \item \textbf{Σωστό}.
    \item \textbf{Λάθος} (π.χ. $f(x) = \cos x$ έχει άπειρες θέσεις ολικού μεγίστου).
\end{alist}

\askhsh
\begin{alist}
    \item Όχι, αφού η $f$ είναι γνησίως αύξουσα (συνεχής συνάρτηση με θετική παράγωγο παντού εκτός ενός σημείου).
    \item Γνησίως αύξουσα στο $\mathbb{R}$.
    \item $f(x) = x^3$.
    \item Αν δεν είναι συνεχής, μπορεί να έχει (π.χ. ασυνέχεια άλματος).
\end{alist}

\askhsh
\begin{alist}
    \item $f(x) = x^3$: $f'(x) = 3x^2 \ge 0$, $f'(0) = 0$ αλλά καμπής, όχι ακρότατο.
    \item $f(x) = x^4$: $f'(x) = 4x^3$, αλλαγή προσήμου, ελάχιστο στο $0$.
    \item $f(x) = |x|$: Ελάχιστο στο $0$ (μη παραγωγίσιμη εκεί).
    \item $f(x) = e^x$: Γνησίως αύξουσα, κανένα ακρότατο.
\end{alist}

\end{document}
